\begin{question}{Limitations of Transformers on Simple Languages}

\begin{subquestion}
    Absorb the value projection $\mathbf{W}_V^{(\ell)}$ into MLP $\mathbf{f}$ by exhibiting a function $\mathbf{f}'$ s.t.
    \begin{equation*}
        \mathbf{x}_t^{(\ell+1)} \;=\; \mathbf{f}'\!\left(\sum_{t'=1}^{T} s_{t,t'}^{(\ell)}\, \mathbf{x}_{t'}^{(\ell)},\; \mathbf{x}_t^{(\ell)}\right) \;=\; \mathbf{f}\!\left(\sum_{t'=1}^{T} s_{t,t'}^{(\ell)}\, \mathbf{v}_{t'}^{(\ell)},\; \mathbf{x}_t^{(\ell)}\right)
    \end{equation*}

    \paragraph{Express $\mathbf{v}_{t'}^{(\ell)}$ in terms of $\mathbf{x}_{t'}^{(\ell)}$}
    From Eq.~(39c), $\mathbf{V}^{(\ell)} = \mathbf{X}^{(\ell)} \mathbf{W}_V^{(\ell)}$. Reading off the $t'$-th row (where the rows of $\mathbf{V}^{(\ell)}$ and $\mathbf{X}^{(\ell)}$ are $\mathbf{v}_{t'}^{(\ell)\top}$ and $\mathbf{x}_{t'}^{(\ell)\top}$ respectively):
    \begin{equation*}
        \mathbf{v}_{t'}^{(\ell)\top} \;=\; \mathbf{x}_{t'}^{(\ell)\top}\, \mathbf{W}_V^{(\ell)}
        \quad\Longleftrightarrow\quad
        \mathbf{v}_{t'}^{(\ell)} \;=\; \mathbf{W}_V^{(\ell)\top}\, \mathbf{x}_{t'}^{(\ell)}
    \end{equation*}

    \paragraph{Pull $\mathbf{W}_V^{(\ell)\top}$ out of the attention sum}
    Substituting into Eq.~(39e), and using that $s_{t,t'}^{(\ell)} \in \mathbb{R}$ is a scalar and matrix multiplication is linear,
    \begin{equation*}
        \mathbf{z}_t^{(\ell)} \;=\; \sum_{t'=1}^{T} s_{t,t'}^{(\ell)}\, \mathbf{v}_{t'}^{(\ell)}
        \;=\; \sum_{t'=1}^{T} s_{t,t'}^{(\ell)}\, \mathbf{W}_V^{(\ell)\top}\, \mathbf{x}_{t'}^{(\ell)}
        \;=\; \mathbf{W}_V^{(\ell)\top} \!\left(\sum_{t'=1}^{T} s_{t,t'}^{(\ell)}\, \mathbf{x}_{t'}^{(\ell)}\right)
    \end{equation*}

    \paragraph{Define $\mathbf{f}'$}
    Plugging this into Eq.~(39f), $\mathbf{x}_t^{(\ell+1)} = \mathbf{f}\!\left(\mathbf{W}_V^{(\ell)\top} \sum_{t'} s_{t,t'}^{(\ell)} \mathbf{x}_{t'}^{(\ell)},\; \mathbf{x}_t^{(\ell)}\right)$
    Hence
    \begin{equation*}
        \mathbf{f}'(\mathbf{a}, \mathbf{b}) \;\defeq\; \mathbf{f}\!\left(\mathbf{W}_V^{(\ell)\top}\, \mathbf{a},\; \mathbf{b}\right) \qquad \text{for } \mathbf{a}, \mathbf{b} \in \mathbb{R}^{D}
    \end{equation*}
    satisfies Eq.~(44), absorbing the linear projection $\mathbf{W}_V^{(\ell)}$ into the first argument of the MLP.

    \paragraph{Lipschitz constant of $\mathbf{f}'$}
    Since $\mathbf{f}$ is $L_{\mathbf{f}}$-Lipschitz and the linear map $(\mathbf{a}, \mathbf{b}) \mapsto (\mathbf{W}_V^{(\ell)\top}\mathbf{a}, \mathbf{b})$ is Lipschitz with constant $\le \sqrt{\|\mathbf{W}_V^{(\ell)}\|_{2}^{2} + 1}$ in the product norm, the composition of Lipschitz functions is Lipschitz. 
    Hence $\mathbf{f}'$ is Lipschitz with constant
    \begin{equation*}
        L_{\mathbf{f}'} \;\le\; L_{\mathbf{f}}\,\sqrt{\|\mathbf{W}_V^{(\ell)}\|_{2}^{2} + 1},
    \end{equation*}
    a finite quantity depending only on $L_{\mathbf{f}}$ and the operator norm of the value projection. In the sequel we abuse notation and write $L_{\mathbf{f}}$ for $L_{\mathbf{f}'}$.
\end{subquestion}

\begin{subquestion}
    Prove $\|\mathbf{x}_t^{(\ell)}\| \le (2L_{\mathbf{f}})^{\ell}\,\|\mathbf{X}^{(0)}\|_{2,\infty}$ by induction on $\ell$. 
    After part~(a), take the layer update as
    \begin{equation*}
        \mathbf{x}_t^{(\ell+1)} \;=\; \mathbf{f}\!\left(\mathbf{z}_t^{(\ell)},\; \mathbf{x}_t^{(\ell)}\right), \qquad \mathbf{z}_t^{(\ell)} \;=\; \sum_{t'=1}^{T} s_{t,t'}^{(\ell)}\, \mathbf{x}_{t'}^{(\ell)}
    \end{equation*}

    \paragraph{Base case ($\ell = 0$)}
    By Definition~10, $\|\mathbf{X}^{(0)}\|_{2,\infty} = \max_{t'} \|\mathbf{x}_{t'}^{(0)}\|$, so for every $t$
    \begin{equation*}
        \|\mathbf{x}_t^{(0)}\| \;\le\; \|\mathbf{X}^{(0)}\|_{2,\infty} \;=\; (2L_{\mathbf{f}})^{0}\,\|\mathbf{X}^{(0)}\|_{2,\infty}
    \end{equation*}

    \paragraph{Inductive hypothesis}
    Assume that for some $\ell \ge 0$ and every $t' \in \{1, \dots, T\}$,
    \begin{equation*}
        \|\mathbf{x}_{t'}^{(\ell)}\| \;\le\; (2L_{\mathbf{f}})^{\ell}\,\|\mathbf{X}^{(0)}\|_{2,\infty}
    \end{equation*}

    \paragraph{Inductive step ($\ell \to \ell+1$)}
    \emph{(i) Bound on $\|\mathbf{z}_t^{(\ell)}\|$} 
    The weights $s_{t,t'}^{(\ell)}$ are softmax outputs, hence $s_{t,t'}^{(\ell)} \ge 0$ and $\sum_{t'=1}^{T} s_{t,t'}^{(\ell)} = 1$
    By the triangle inequality and the inductive hypothesis,
    \begin{equation*}
        \|\mathbf{z}_t^{(\ell)}\| \;=\; \Bigg\|\sum_{t'=1}^{T} s_{t,t'}^{(\ell)}\, \mathbf{x}_{t'}^{(\ell)}\Bigg\| \;\le\; \sum_{t'=1}^{T} s_{t,t'}^{(\ell)}\, \|\mathbf{x}_{t'}^{(\ell)}\| \;\le\; (2L_{\mathbf{f}})^{\ell}\,\|\mathbf{X}^{(0)}\|_{2,\infty}
    \end{equation*}

    \emph{(ii) Lipschitz step} Using $\mathbf{f}(\mathbf{0}, \mathbf{0}) = \mathbf{0}$ and the Lipschitz property (Definition~9) of $\mathbf{f}: \mathbb{R}^{D} \times \mathbb{R}^{D} \to \mathbb{R}^{D}$ with the product space norm $\|(\mathbf{a}, \mathbf{b})\| = \sqrt{\|\mathbf{a}\|^{2} + \|\mathbf{b}\|^{2}}$,
    \begin{equation*}
        \|\mathbf{f}(\mathbf{z}, \mathbf{x})\| \;=\; \|\mathbf{f}(\mathbf{z}, \mathbf{x}) - \mathbf{f}(\mathbf{0}, \mathbf{0})\| \;\le\; L_{\mathbf{f}}\sqrt{\|\mathbf{z}\|^{2} + \|\mathbf{x}\|^{2}} \;\le\; L_{\mathbf{f}}\bigl(\|\mathbf{z}\| + \|\mathbf{x}\|\bigr),
    \end{equation*}
    where the last inequality follows from $\|\mathbf{z}\|^{2} + \|\mathbf{x}\|^{2} \le \bigl(\|\mathbf{z}\| + \|\mathbf{x}\|\bigr)^{2}$.

    \emph{(iii)} Applying this to $\mathbf{x}_t^{(\ell+1)} = \mathbf{f}(\mathbf{z}_t^{(\ell)}, \mathbf{x}_t^{(\ell)})$ and substituting both bounds,
    \begin{align*}
        \|\mathbf{x}_t^{(\ell+1)}\|
        &\;\le\; L_{\mathbf{f}}\bigl(\|\mathbf{z}_t^{(\ell)}\| + \|\mathbf{x}_t^{(\ell)}\|\bigr) \\
        &\;\le\; L_{\mathbf{f}}\bigl((2L_{\mathbf{f}})^{\ell}\,\|\mathbf{X}^{(0)}\|_{2,\infty} + (2L_{\mathbf{f}})^{\ell}\,\|\mathbf{X}^{(0)}\|_{2,\infty}\bigr) \\
        &\;=\; (2L_{\mathbf{f}})^{\ell+1}\,\|\mathbf{X}^{(0)}\|_{2,\infty}
    \end{align*}
    By induction, $\|\mathbf{x}_t^{(\ell)}\| \le (2L_{\mathbf{f}})^{\ell}\,\|\mathbf{X}^{(0)}\|_{2,\infty}$ for every $\ell \ge 0$ and every $t \in \{1, \dots, T\}$
\end{subquestion}

\begin{subquestion}
    Fix a layer $\ell$ and a query position $t$. The pre softmax logit vector is $\mathbf{l} = \mathbf{K}^{(\ell)} \mathbf{q}_t^{(\ell)} \in \mathbb{R}^{T}$, with $t'$-th entry
    \begin{equation*}
        l_{t'} \;=\; \mathbf{k}_{t'}^{(\ell)\top}\, \mathbf{q}_{t}^{(\ell)}
    \end{equation*}

    \paragraph{Rewrite the logit in terms of $\mathbf{x}^{(\ell)}$}
    From Eqs.~(39a) and (39b), reading off rows as in part~(a) gives $\mathbf{q}_t^{(\ell)} = \mathbf{W}_Q^{(\ell)\top} \mathbf{x}_t^{(\ell)}$ and $\mathbf{k}_{t'}^{(\ell)} = \mathbf{W}_K^{(\ell)\top} \mathbf{x}_{t'}^{(\ell)}$. Substituting,
    \begin{equation*}
        l_{t'} \;=\; \bigl(\mathbf{W}_K^{(\ell)\top}\, \mathbf{x}_{t'}^{(\ell)}\bigr)^{\!\top}\bigl(\mathbf{W}_Q^{(\ell)\top}\, \mathbf{x}_{t}^{(\ell)}\bigr) \;=\; \mathbf{x}_{t'}^{(\ell)\top}\, \mathbf{W}_K^{(\ell)}\, \mathbf{W}_Q^{(\ell)\top}\, \mathbf{x}_{t}^{(\ell)}
    \end{equation*}

    \paragraph{Apply Cauchy-Schwarz and the operator norm inequality}
    By Cauchy-Schwarz and $\|\mathbf{M}\mathbf{v}\| \le \|\mathbf{M}\|_{2}\,\|\mathbf{v}\|$ for the induced 2-norm,
    \begin{equation*}
        |l_{t'}| \;\le\; \|\mathbf{x}_{t'}^{(\ell)}\| \cdot \|\mathbf{W}_K^{(\ell)}\, \mathbf{W}_Q^{(\ell)\top}\, \mathbf{x}_{t}^{(\ell)}\| \;\le\; \|\mathbf{x}_{t'}^{(\ell)}\| \cdot \|\mathbf{W}_K^{(\ell)}\, \mathbf{W}_Q^{(\ell)\top}\|_{2} \cdot \|\mathbf{x}_{t}^{(\ell)}\|
    \end{equation*}

    \paragraph{Bound the representation norms uniformly by $F$}
    From part~(b), $\|\mathbf{x}_t^{(\ell)}\| \le (2L_{\mathbf{f}})^{\ell}\,\|\mathbf{X}^{(0)}\|_{2,\infty}$. With $F \ge \max\!\bigl(1,\, (2L_{\mathbf{f}})^{L}\bigr)\,\|\mathbf{X}^{(0)}\|_{2,\infty}$, we have for every $\ell \in \{0, \dots, L\}$:
    \begin{itemize}
        \item if $2L_{\mathbf{f}} \ge 1$: $(2L_{\mathbf{f}})^{\ell} \le (2L_{\mathbf{f}})^{L}$, so $\|\mathbf{x}_t^{(\ell)}\| \le (2L_{\mathbf{f}})^{L}\,\|\mathbf{X}^{(0)}\|_{2,\infty} \le F$
        \item if $2L_{\mathbf{f}} < 1$: $(2L_{\mathbf{f}})^{\ell} \le 1$, so $\|\mathbf{x}_t^{(\ell)}\| \le \|\mathbf{X}^{(0)}\|_{2,\infty} \le F$
    \end{itemize}
    Hence $\|\mathbf{x}_t^{(\ell)}\|, \|\mathbf{x}_{t'}^{(\ell)}\| \le F$, and
    \begin{equation*}
        |l_{t'}| \;\le\; F^{2}\, \|\mathbf{W}_K^{(\ell)}\, \mathbf{W}_Q^{(\ell)\top}\|_{2}
    \end{equation*}

    \paragraph{Define $L_{\mathrm{att}}$}
    Since the parameter matrices are fixed, the quantity
    \begin{equation*}
        L_{\mathrm{att}} \;\defeq\; \max_{\ell \in \{1, \dots, L\}}\, \|\mathbf{W}_K^{(\ell)}\, \mathbf{W}_Q^{(\ell)\top}\|_{2}
    \end{equation*}
    is a finite constant that depends only on the model parameters. Setting $A \defeq F^{2} L_{\mathrm{att}}$, the bound $|l_{t'}| \le A$ holds for every entry $t' \in \{1, \dots, T\}$, hence
    \begin{equation*}
        \|\mathbf{l}\|_{\infty} \;=\; \max_{t'} |l_{t'}| \;\le\; F^{2}\, L_{\mathrm{att}} \;=\; A
    \end{equation*}
\end{subquestion}

\begin{subquestion}
    Let $\mathbf{l} \in \mathbb{R}^{T}$ with $\|\mathbf{l}\|_{\infty} \le A$ and write $\mathbf{s} = \operatorname{softmax}(\mathbf{l})$, i.e., $s_t = \exp(l_t) / \sum_{t'=1}^{T} \exp(l_{t'})$.

    \paragraph{Per coordinate bounds on $\exp(l_t)$}
    Since $\|\mathbf{l}\|_{\infty} \le A$, every coordinate satisfies $-A \le l_t \le A$. 
    Because $\exp$ is monotonically increasing,
    \begin{equation*}
        \exp(-A) \;\le\; \exp(l_t) \;\le\; \exp(A) \qquad \text{for all } t \in \{1, \dots, T\}
    \end{equation*}

    Summing the per coordinate bounds over $t' \in \{1, \dots, T\}$ for the softmax denominator,
    \begin{equation*}
        T\,\exp(-A) \;\le\; \sum_{t'=1}^{T} \exp(l_{t'}) \;\le\; T\,\exp(A)
    \end{equation*}

    \paragraph{Combine for $s_t$}
    The fraction $s_t = \exp(l_t) / \sum_{t'} \exp(l_{t'})$ is maximized by pairing the largest numerator with the smallest denominator, and minimized by the reverse. 
    Using the bounds,
    \begin{equation*}
        s_t \;\le\; \frac{\exp(A)}{T\,\exp(-A)} \;=\; \frac{\exp(2A)}{T},
        \qquad
        s_t \;\ge\; \frac{\exp(-A)}{T\,\exp(A)} \;=\; \frac{\exp(-2A)}{T}
    \end{equation*}
    Hence
    \begin{equation*}
        \frac{\exp(-2A)}{T} \;\le\; s_t \;\le\; \frac{\exp(2A)}{T} \qquad \text{for every } t \in \{1, \dots, T\},
    \end{equation*}
\end{subquestion}

\begin{subquestion}
    Using the simplified attention from part~(a) that and $\Delta_t^{(\ell)} = \|\mathbf{z}_t^{(\ell)} - \mathbf{z}_t^{\prime(\ell)}\|$ as in Eq.~(48).
    \begin{equation*}
        \mathbf{z}_t^{(\ell)} \;=\; \sum_{t'=1}^{T} s_{t,t'}^{(\ell)}\, \mathbf{x}_{t'}^{(\ell)},
        \qquad
        \mathbf{z}_t^{\prime(\ell)} \;=\; \sum_{t'=1}^{T} s_{t,t'}^{\prime(\ell)}\, \mathbf{x}_{t'}^{\prime(\ell)},
    \end{equation*}
    \paragraph{Add and subtract a cross term}
    Insert $\pm \sum_{t'} s_{t,t'}^{\prime(\ell)}\, \mathbf{x}_{t'}^{(\ell)}$ inside the difference and group:
    \begin{align*}
        \mathbf{z}_t^{(\ell)} - \mathbf{z}_t^{\prime(\ell)}
        &\;=\; \sum_{t'=1}^{T} s_{t,t'}^{(\ell)}\, \mathbf{x}_{t'}^{(\ell)} \;-\; \sum_{t'=1}^{T} s_{t,t'}^{\prime(\ell)}\, \mathbf{x}_{t'}^{\prime(\ell)} \\
        &\;=\; \sum_{t'=1}^{T} \bigl(s_{t,t'}^{(\ell)} - s_{t,t'}^{\prime(\ell)}\bigr)\, \mathbf{x}_{t'}^{(\ell)}
        \;+\; \sum_{t'=1}^{T} s_{t,t'}^{\prime(\ell)}\, \bigl(\mathbf{x}_{t'}^{(\ell)} - \mathbf{x}_{t'}^{\prime(\ell)}\bigr)
    \end{align*}
    The first sum isolates the change in attention weights with the original representations, the second isolates the change in representations weighted by the perturbed attention.

    Taking norms, applying the triangle inequality on each outer sum and on each scalar vector product, and using $s_{t,t'}^{\prime(\ell)} \ge 0$ (so $|s_{t,t'}^{\prime(\ell)}| = s_{t,t'}^{\prime(\ell)}$) gives Eq.~(49):
    \begin{equation*}
        \Delta_t^{(\ell)} \;\le\; \underbrace{\sum_{t'=1}^{T} \bigl|s_{t,t'}^{(\ell)} - s_{t,t'}^{\prime(\ell)}\bigr|\, \|\mathbf{x}_{t'}^{(\ell)}\|}_{\text{Difficult Part}}
        \;+\; \underbrace{\sum_{t'=1}^{T} s_{t,t'}^{\prime(\ell)}\, \|\mathbf{x}_{t'}^{(\ell)} - \mathbf{x}_{t'}^{\prime(\ell)}\|}_{\text{Easy Part}}
    \end{equation*}
\end{subquestion}

\begin{subquestion}
    The Difficult Part (Eq.~(50)) is $\sum_{t'=1}^{T} \bigl|s_{t,t'}^{(\ell)} - s_{t,t'}^{\prime(\ell)}\bigr|\, \|\mathbf{x}_{t'}^{(\ell)}\|$. 
    We have a uniform bound on each $\|\mathbf{x}_{t'}^{(\ell)}\|$ and definition~10 gives $\|\mathbf{X}^{(\ell)}\|_{2,\infty} = \max_{t''} \|\mathbf{x}_{t''}^{(\ell)}\|$, so for every $t' \in \{1, \dots, T\}$,
    \begin{equation*}
        \|\mathbf{x}_{t'}^{(\ell)}\| \;\le\; \|\mathbf{X}^{(\ell)}\|_{2,\infty}
    \end{equation*}

    Since $\bigl|s_{t,t'}^{(\ell)} - s_{t,t'}^{\prime(\ell)}\bigr| \ge 0$, multiplying through preserves the inequality termwise. 
    Summing over $t'$ and pulling the constant $\|\mathbf{X}^{(\ell)}\|_{2,\infty}$ out of the sum yields Eq.~(51):
    \begin{equation*}
        \sum_{t'=1}^{T} \bigl|s_{t,t'}^{(\ell)} - s_{t,t'}^{\prime(\ell)}\bigr|\, \|\mathbf{x}_{t'}^{(\ell)}\|
        \;\le\; \sum_{t'=1}^{T} \bigl|s_{t,t'}^{(\ell)} - s_{t,t'}^{\prime(\ell)}\bigr|\, \|\mathbf{X}^{(\ell)}\|_{2,\infty}
        \;=\; \!\left(\sum_{t'=1}^{T} \bigl|s_{t,t'}^{(\ell)} - s_{t,t'}^{\prime(\ell)}\bigr|\right)\!\|\mathbf{X}^{(\ell)}\|_{2,\infty}
    \end{equation*}
\end{subquestion}

\begin{subquestion}
    First prove the general claim and then apply it to the exponential. Let $f: \mathbb{R} \to \mathbb{R}$ be twice differentiable, and let $(a_T)_{T \in \mathbb{N}}$, $(b_T)_{T \in \mathbb{N}}$ be bounded with $|a_T - b_T| = \mathcal{O}(1/T)$. We claim $|f(a_T) - f(b_T)| = \mathcal{O}(1/T)$.

    \emph{(i) Uniform bound on $|f'|$} Since both sequences are bounded, fix $M > 0$ with $a_T, b_T \in [-M, M]$ for every $T$. Twice differentiability of $f$ implies $f'$ is differentiable, hence continuous. By the Extreme Value Theorem, $f'$ attains its maximum on the compact interval $[-M, M]$:
    \begin{equation*}
        K \;\defeq\; \max_{x \in [-M, M]} |f'(x)| \;<\; \infty
    \end{equation*}

    \emph{(ii) Mean value theorem} For each $T$, the MVT gives some $c_T$ between $a_T$ and $b_T$ (so $c_T \in [-M, M]$) with $f(a_T) - f(b_T) = f'(c_T)\,(a_T - b_T)$. Hence
    \begin{equation*}
        |f(a_T) - f(b_T)| \;=\; |f'(c_T)|\, |a_T - b_T| \;\le\; K\, |a_T - b_T|
    \end{equation*}

    \emph{(iii) Asymptotic conclusion} Since $|a_T - b_T| = \mathcal{O}(1/T)$ and multiplication by the constant $K$ preserves the order,
    \begin{equation*}
        |f(a_T) - f(b_T)| \;\le\; K \cdot \mathcal{O}(1/T) \;=\; \mathcal{O}(1/T)
    \end{equation*}

    \paragraph{Application to $f = \exp$}
    The function $\exp: \mathbb{R} \to \mathbb{R}$ is $C^{\infty}$, hence twice differentiable. Applying the result with $f = \exp$ to bounded $(a_T), (b_T)$ with $|a_T - b_T| = \mathcal{O}(1/T)$ gives
    \begin{equation*}
        |\exp(a_T) - \exp(b_T)| \;=\; \mathcal{O}(1/T)
    \end{equation*}
\end{subquestion}

\begin{subquestion}
    Let $\mathbf{s}(\mathbf{x}) = \operatorname{softmax}(\mathbf{x})$ and write $s_t(\mathbf{x})$ for its $t$-th coordinate. 
    The boundedness assumption gives some $A > 0$ with $\|\mathbf{a}_T\|_{\infty}, \|\mathbf{b}_T\|_{\infty} \le A$ for every $T$.

    \paragraph{Mean value theorem on $s_t$}
    Since $s_t$ is $C^{\infty}$ on $\mathbb{R}^{T}$, the multivariate MVT gives a point $\mathbf{c}_T = \lambda \mathbf{a}_T + (1-\lambda)\mathbf{b}_T$ on the segment between $\mathbf{a}_T$ and $\mathbf{b}_T$ s.t.
    \begin{equation*}
        s_t(\mathbf{a}_T) - s_t(\mathbf{b}_T) \;=\; \nabla s_t(\mathbf{c}_T)^{\!\top}\bigl(\mathbf{a}_T - \mathbf{b}_T\bigr) \;=\; \sum_{k=1}^{T} \frac{\partial s_t}{\partial x_k}(\mathbf{c}_T)\,\bigl((\mathbf{a}_T)_k - (\mathbf{b}_T)_k\bigr)
    \end{equation*}
    Convexity of $\|\cdot\|_{\infty}$ also gives $\|\mathbf{c}_T\|_{\infty} \le A$, so by part~(d),
    \begin{equation*}
        s_k(\mathbf{c}_T) \;\le\; \frac{\exp(2A)}{T} \;=\; \mathcal{O}(1/T) \qquad \text{for every } k
    \end{equation*}

    \paragraph{Softmax Jacobian}
    The standard formula is $\partial s_t / \partial x_k = s_t\,\delta_{tk} - s_t s_k$, i.e., $s_t(1-s_t)$ on the diagonal and $-s_t s_k$ off diagonal. 
    Splitting the sum at $k = t$,
    \begin{equation*}
        s_t(\mathbf{a}_T) - s_t(\mathbf{b}_T)
        \;=\; \underbrace{s_t(\mathbf{c}_T)\bigl(1 - s_t(\mathbf{c}_T)\bigr)\bigl((\mathbf{a}_T)_t - (\mathbf{b}_T)_t\bigr)}_{\text{Term 1}}
        \;-\; \underbrace{\sum_{k \ne t} s_t(\mathbf{c}_T)\, s_k(\mathbf{c}_T)\,\bigl((\mathbf{a}_T)_k - (\mathbf{b}_T)_k\bigr)}_{\text{Term 2}}
    \end{equation*}

    Using $0 \le 1 - s_t(\mathbf{c}_T) \le 1$, the part~(d) bound on $s_t(\mathbf{c}_T)$, and the hypothesis on the difference,
    \begin{equation*}
        |\text{Term 1}| \;\le\; s_t(\mathbf{c}_T)\,\bigl|(\mathbf{a}_T)_t - (\mathbf{b}_T)_t\bigr|
        \;=\; \mathcal{O}(1/T) \cdot \mathcal{O}(1/T) \;=\; \mathcal{O}(1/T^{2})
    \end{equation*}

    Pull $s_t(\mathbf{c}_T)$ outside the sum, bound the differences uniformly by $\max_k |(\mathbf{a}_T)_k - (\mathbf{b}_T)_k| = \mathcal{O}(1/T)$, and use $\sum_{k \ne t} s_k(\mathbf{c}_T) \le \sum_k s_k(\mathbf{c}_T) = 1$:
    \begin{equation*}
        |\text{Term 2}| \;\le\; s_t(\mathbf{c}_T)\,\Bigl(\max_k \bigl|(\mathbf{a}_T)_k - (\mathbf{b}_T)_k\bigr|\Bigr)\sum_{k \ne t} s_k(\mathbf{c}_T)
        \;\le\; \mathcal{O}(1/T) \cdot \mathcal{O}(1/T) \cdot 1 \;=\; \mathcal{O}(1/T^{2})
    \end{equation*}

    By the triangle inequality the following equals to Eq.~(52):
    \begin{equation*}
        |\operatorname{softmax}(\mathbf{a}_T)_t - \operatorname{softmax}(\mathbf{b}_T)_t| \;\le\; |\text{Term 1}| + |\text{Term 2}| \;=\; \mathcal{O}(1/T^{2}),
    \end{equation*}
\end{subquestion}

\begin{subquestion}
    Let's translate the hypothesis to vector norms.
    The assumption $\|\mathbf{w}_{T,d}^{(\ell)}\|_{\infty} = \mathcal{O}(1/T)$ for every $d \in [D]$ means $|(\mathbf{x}_{t'}^{(\ell)} - \mathbf{x}_{t'}^{\prime(\ell)})_d| = \mathcal{O}(1/T)$ for every $d \in [D]$ and every $t' \ne t^*$. 
    Since $D$ is a model constant, summing the $D$ squared bounds and taking the square root gives:
    \begin{equation*}
        \|\mathbf{x}_{t'}^{(\ell)} - \mathbf{x}_{t'}^{\prime(\ell)}\| \;\le\; \sqrt{D}\, \max_d \|\mathbf{w}_{T,d}^{(\ell)}\|_{\infty} \;=\; \mathcal{O}(1/T) \qquad \text{for every } t' \ne t^*
    \end{equation*}

    \paragraph{Bound the softmax differences via part~(h)}
    Throughout this paragraph we assume the query position satisfies $t \ne t^*$, so $\|\mathbf{x}_t^{(\ell)} - \mathbf{x}_t^{\prime(\ell)}\| = \mathcal{O}(1/T)$ on the query side too.
    The $t = t^*$ case is handled by a trivial bound in part~(j).
    The attention logits $\mathbf{l}_t^{(\ell)} = \mathbf{K}^{(\ell)} \mathbf{q}_t^{(\ell)}$ are bounded by part~(c).
    Splitting the bilinear logit difference and using part~(b) to bound the remaining factors by $\mathcal{O}(1)$,
    \begin{equation*}
        l_{t,t'}^{(\ell)} - l_{t,t'}^{\prime(\ell)} \;=\; (\mathbf{x}_t^{(\ell)} - \mathbf{x}_t^{\prime(\ell)})^{\!\top}\mathbf{W}_K^{(\ell)}\mathbf{W}_Q^{(\ell)\top}\mathbf{x}_{t'}^{(\ell)} \;+\; \mathbf{x}_t^{\prime(\ell)\top}\mathbf{W}_K^{(\ell)}\mathbf{W}_Q^{(\ell)\top}(\mathbf{x}_{t'}^{(\ell)} - \mathbf{x}_{t'}^{\prime(\ell)}),
    \end{equation*}
    both differences are $\mathcal{O}(1/T)$ when $t, t' \ne t^*$, hence $|l_{t,t'}^{(\ell)} - l_{t,t'}^{\prime(\ell)}| = \mathcal{O}(1/T)$ for $t' \ne t^*$.
    Applying part~(h) to the resulting softmax outputs yields
    \begin{equation*}
        \bigl|s_{t,t'}^{(\ell)} - s_{t,t'}^{\prime(\ell)}\bigr| \;=\; \mathcal{O}(1/T^{2}) \qquad \text{for every } t' \ne t^*
    \end{equation*}

    \paragraph{Bound the representations via part~(b)}
    Part~(b) gives $\|\mathbf{x}_{t'}^{(\ell)}\| \le (2L_{\mathbf{f}})^{\ell}\,\|\mathbf{X}^{(0)}\|_{2,\infty}$. 
    Since $L$ is fixed, this is $\mathcal{O}(1)$ uniformly in $T$ and $t'$.

    Substituting the two bounds into the truncated Difficult Part and using that the sum has $T - 1$ terms, we get Eq.~(53):
    \begin{equation*}
        \sum_{\substack{t'=1 \\ t' \ne t^*}}^{T} \bigl|s_{t,t'}^{(\ell)} - s_{t,t'}^{\prime(\ell)}\bigr|\, \|\mathbf{x}_{t'}^{(\ell)}\|
        \;\le\; (T-1) \cdot \mathcal{O}(1/T^{2}) \cdot \mathcal{O}(1)
        \;=\; \mathcal{O}(1/T),
    \end{equation*}
\end{subquestion}

\begin{subquestion}
    Prove Eqs.~(55-56) jointly by induction on $\ell$, then apply the bound to $\mathbf{h}(\mathbf{y})$.

    \paragraph{Base case ($\ell = 0$)}
    By Eq.~(38), $\mathbf{x}_t^{(0)} = (\mathbf{r}(y_t),\, \mathbf{u}(t))$. Since $\mathbf{y}$ and $\mathbf{y}'$ differ only at $t^*$:
    \begin{itemize}
        \item if $t \ne t^*$: $\mathbf{x}_t^{(0)} = \mathbf{x}_t^{\prime(0)}$, so $\|\mathbf{x}_t^{(0)} - \mathbf{x}_t^{\prime(0)}\| = 0 = \mathcal{O}(1/T)$
        \item if $t = t^*$: $\|\mathbf{x}_{t^*}^{(0)} - \mathbf{x}_{t^*}^{\prime(0)}\| = \delta = \mathcal{O}(1)$ (a constant in $T$, since $\mathbf{r}$ ranges over finitely many bounded embeddings).
    \end{itemize}
    So Eq.~(56) holds at $\ell = 0$.

    \paragraph{Inductive hypothesis}
    Assume Eq.~(56) holds at layer $\ell$:
    \begin{equation*}
        \|\mathbf{x}_t^{(\ell)} - \mathbf{x}_t^{\prime(\ell)}\| \;=\;
        \begin{cases} \mathcal{O}(1) & \text{if } t = t^*, \\ \mathcal{O}(1/T) & \text{otherwise.} \end{cases}
    \end{equation*}

    \paragraph{Inductive step: bounding $\Delta_t^{(\ell)}$, Eq.~(55)}
    The bound depends on whether the query position $t$ is the flipped one, as the bilinear logit propagation in part~(i) requires the query side to also differ by $\mathcal{O}(1/T)$. 
    Split into two cases.

    \emph{Case $t \ne t^*$} Take Eq.~(54b)'s four terms and bound each:
    \begin{itemize}
        \item \emph{Term 1} ($\sum_{t' \ne t^*} |s - s'|\,\|\mathbf{x}_{t'}^{(\ell)}\|$): the inductive hypothesis says $\|\mathbf{x}_{t'}^{(\ell)} - \mathbf{x}_{t'}^{\prime(\ell)}\| = \mathcal{O}(1/T)$ for $t' \ne t^*$ and for the query $t$ itself (since $t \ne t^*$), so part~(i) applies and gives $\mathcal{O}(1/T)$.
        \item \emph{Term 2} ($|s_{t,t^*}^{(\ell)} - s_{t,t^*}^{\prime(\ell)}|\,\|\mathbf{x}_{t^*}^{(\ell)}\|$): part~(d) gives $|s_{t,t^*}^{(\ell)}|, |s_{t,t^*}^{\prime(\ell)}| \le \exp(2A)/T$, so the difference is $\mathcal{O}(1/T)$, and part~(b) gives $\|\mathbf{x}_{t^*}^{(\ell)}\| = \mathcal{O}(1)$. Product: $\mathcal{O}(1/T)$.
        \item \emph{Term 3} ($\sum_{t' \ne t^*} s_{t,t'}^{\prime(\ell)}\,\|\mathbf{x}_{t'}^{(\ell)} - \mathbf{x}_{t'}^{\prime(\ell)}\|$): the inductive hypothesis bounds each $\|\mathbf{x}_{t'}^{(\ell)} - \mathbf{x}_{t'}^{\prime(\ell)}\| = \mathcal{O}(1/T)$, so we can pull it out and use $\sum_{t' \ne t^*} s_{t,t'}^{\prime(\ell)} \le 1$. Product: $\mathcal{O}(1/T)$.
        \item \emph{Term 4} ($s_{t,t^*}^{\prime(\ell)}\,\|\mathbf{x}_{t^*}^{(\ell)} - \mathbf{x}_{t^*}^{\prime(\ell)}\|$): part~(d) gives $s_{t,t^*}^{\prime(\ell)} = \mathcal{O}(1/T)$ and the inductive hypothesis gives $\|\mathbf{x}_{t^*}^{(\ell)} - \mathbf{x}_{t^*}^{\prime(\ell)}\| = \mathcal{O}(1)$. Product: $\mathcal{O}(1/T)$.
    \end{itemize}
    Summing the four $\mathcal{O}(1/T)$ contributions yields $\Delta_t^{(\ell)} = \mathcal{O}(1/T)$ for $t \ne t^*$.

    \emph{Case $t = t^*$} 
    The query side $\mathbf{x}_{t^*}^{(\ell)} - \mathbf{x}_{t^*}^{\prime(\ell)}$ is only $\mathcal{O}(1)$ by the inductive hypothesis, so the bilinear logit propagation in part~(i) does not apply at the query position and Term~1 cannot be sharpened. 
    Using the trivial bound $\sum_{t' \ne t^*} |s_{t^*,t'}^{(\ell)} - s_{t^*,t'}^{\prime(\ell)}| \le \sum_{t' \ne t^*} (s_{t^*,t'}^{(\ell)} + s_{t^*,t'}^{\prime(\ell)}) \le 2$ together with $\|\mathbf{x}_{t'}^{(\ell)}\| \le F = \mathcal{O}(1)$ from part~(b) gives Term~1 $= \mathcal{O}(1)$. 
    Terms~2-4 are still $\mathcal{O}(1/T)$ by the same arguments above as none of them invokes part~(h) on the query side.
    Hence $\Delta_{t^*}^{(\ell)} = \mathcal{O}(1)$.

    Combining the two cases gives Eq.~(55) in the form
    \begin{equation*}
        \Delta_t^{(\ell)} \;=\;
        \begin{cases} \mathcal{O}(1) & \text{if } t = t^*, \\ \mathcal{O}(1/T) & \text{otherwise.} \end{cases}
    \end{equation*}

    \paragraph{Inductive step: bounding $\|\mathbf{x}_t^{(\ell+1)} - \mathbf{x}_t^{\prime(\ell+1)}\|$ (Eq.~(56))}
    By the Lipschitz step from part~(b) applied to differences (rather than to $\mathbf{f}(\mathbf{0},\mathbf{0})$),
    \begin{equation*}
        \|\mathbf{x}_t^{(\ell+1)} - \mathbf{x}_t^{\prime(\ell+1)}\|
        \;=\; \|\mathbf{f}(\mathbf{z}_t^{(\ell)}, \mathbf{x}_t^{(\ell)}) - \mathbf{f}(\mathbf{z}_t^{\prime(\ell)}, \mathbf{x}_t^{\prime(\ell)})\|
        \;\le\; L_{\mathbf{f}}\bigl(\Delta_t^{(\ell)} + \|\mathbf{x}_t^{(\ell)} - \mathbf{x}_t^{\prime(\ell)}\|\bigr)
    \end{equation*}
    Splitting on the value of $t$ and combining the two case bound on $\Delta_t^{(\ell)}$ with the inductive hypothesis,
    \begin{equation*}
        \|\mathbf{x}_t^{(\ell+1)} - \mathbf{x}_t^{\prime(\ell+1)}\|
        \;\le\;
        \begin{cases} L_{\mathbf{f}}\bigl(\mathcal{O}(1) + \mathcal{O}(1)\bigr) = \mathcal{O}(1) & \text{if } t = t^*, \\ L_{\mathbf{f}}\bigl(\mathcal{O}(1/T) + \mathcal{O}(1/T)\bigr) = \mathcal{O}(1/T) & \text{otherwise,} \end{cases}
    \end{equation*}
    closing the induction and giving Eq.~(56) at $\ell+1$.

    \paragraph{Implication for $\mathbf{h}(\mathbf{y})$}
    By definition, $\mathbf{h}(\mathbf{y}) = \mathbf{x}_T^{(L)}$. The hypothesis $t^* < T$ means $T \ne t^*$, so the otherwise case of Eq.~(56) applies:
    \begin{equation*}
        \|\mathbf{h}(\mathbf{y}) - \mathbf{h}(\mathbf{y}')\| \;=\; \|\mathbf{x}_T^{(L)} - \mathbf{x}_T^{\prime(L)}\| \;=\; \mathcal{O}(1/T)
    \end{equation*}
    The final \textsc{eos} representations of a string and any one bit flipped variant become indistinguishable at rate $1/T$ as the sequence length grows, this is Lemma~1.
\end{subquestion}

\begin{subquestion}
    Apply Lemma~1 to the \texttt{Odd-1} acceptor and draw out its consequences.

    \paragraph{(i) Parity of $\mathbf{y}'$}
    Suppose $\mathbf{y} \in \texttt{Odd-1}$, so $\mathbf{y}$ contains $2k+1$ ones for some $k \in \mathbb{N}$. 
    Flipping a single bit at position $t^*$ either turns a $0$ into a $1$ (giving $2k+2$ ones) or a $1$ into a $0$ (giving $2k$ ones). 
    Either way, the count's parity flips from odd to even, so
    \begin{equation*}
        \mathbf{y}' \;\notin\; \texttt{Odd-1}
    \end{equation*}
    The correct labels of $\mathbf{y}$ and $\mathbf{y}'$ are therefore opposite.

    \paragraph{(ii) Implication for $\hat{p}'$}
    From part~(j), $\|\mathbf{h}(\mathbf{y}) - \mathbf{h}(\mathbf{y}')\| = \mathcal{O}(1/T)$. 
    Since the readout vector $\mathbf{w}$ is fixed (independent of $T$), Cauchy-Schwarz gives
    \begin{equation*}
        |\mathbf{h}(\mathbf{y})^{\top} \mathbf{w} - \mathbf{h}(\mathbf{y}')^{\top} \mathbf{w}| \;\le\; \|\mathbf{h}(\mathbf{y}) - \mathbf{h}(\mathbf{y}')\| \cdot \|\mathbf{w}\| \;=\; \mathcal{O}(1/T)
    \end{equation*}
    The two scalar logits $\mathbf{h}(\mathbf{y})^{\top}\mathbf{w}$ and $\mathbf{h}(\mathbf{y}')^{\top}\mathbf{w}$ are bounded sequences in $T$ (since $\|\mathbf{h}\|$ is bounded by part~(b) and $\|\mathbf{w}\|$ is fixed), and they differ by $\mathcal{O}(1/T)$. 
    The sigmoid $\sigma$ is $C^{\infty}$ on $\mathbb{R}$, so part~(g) applies and
    \begin{equation*}
        |\hat{p} - \hat{p}'| \;=\; |\sigma(\mathbf{h}(\mathbf{y})^{\top}\mathbf{w}) - \sigma(\mathbf{h}(\mathbf{y}')^{\top}\mathbf{w})| \;=\; \mathcal{O}(1/T)
    \end{equation*}
    Concretely: if the model is correct on $\mathbf{y}$, i.e., $\hat{p} > 0.5$, then for sufficiently large $T$ the gap $|\hat{p} - \hat{p}'|$ is below the margin $\hat{p} - 0.5$, forcing $\hat{p}' > 0.5$ as well. 
    By Eqs.~(41a) and (41b), the model then accepts $\mathbf{y}'$ too, but $\mathbf{y}' \notin \texttt{Odd-1}$, so the model misclassifies $\mathbf{y}'$.

    \paragraph{(iii) Limit $|\mathbf{y}| \to \infty$}
    Taking $T = |\mathbf{y}| \to \infty$ in the bound from (ii),
    \begin{equation*}
        \lim_{T \to \infty} |\hat{p} - \hat{p}'| \;=\; 0
    \end{equation*}
    The transformer assigns asymptotically same probabilities to $\mathbf{y}$ and to its one bit flipped variant $\mathbf{y}'$ that has opposite parity. 
    Hence soft attention transformers cannot recognize \texttt{Odd-1} for long inputs and any decision boundary is eventually breached by a single bit perturbation.
\end{subquestion}

\begin{subquestion}
    How layer normalization affects the chain of bounds, and what the asymptotic nature of the result actually rules out.

    \paragraph{(i) Adding layer normalization}
    Layer normalization, applied per token, has the form
    \begin{equation*}
        \operatorname{LN}(\mathbf{x}) \;=\; \frac{\mathbf{x} - \mu(\mathbf{x})\,\mathbf{1}}{\sqrt{\sigma^{2}(\mathbf{x}) + \epsilon}} \odot \boldsymbol{\gamma} + \boldsymbol{\beta},
    \end{equation*}
    with $\mu, \sigma^{2}$ the empirical mean and variance of the entries of $\mathbf{x}$, and learnable $\boldsymbol{\gamma}, \boldsymbol{\beta}$. 
    The key question is whether $\operatorname{LN}$ is Lipschitz continuous, since the entire proof in parts~(b)-(j) depended on a finite Lipschitz constant for $\mathbf{f}$.

    \emph{Case $\epsilon > 0$} 
    The denominator $\sqrt{\sigma^{2}(\mathbf{x}) + \epsilon} \ge \sqrt{\epsilon} > 0$ is bounded away from zero, so $\operatorname{LN}$ is smooth with a bounded gradient on $\mathbb{R}^{D}$, hence Lipschitz with some constant $L_{\operatorname{LN}}$. 
    The composition $\mathbf{f} \circ \operatorname{LN}$ or $\operatorname{LN} \circ \mathbf{f}$ are then Lipschitz with constant $L_{\mathbf{f}} \cdot L_{\operatorname{LN}}$. 
    Replacing $L_{\mathbf{f}}$ by this product throughout every bound from (b) onward goes through unchanged and the asymptotic $\mathcal{O}(1/T)$ conclusion of Lemma~1 is preserved.

    \emph{Case $\epsilon = 0$} 
    The denominator $\sigma(\mathbf{x})$ vanishes whenever the entries of $\mathbf{x}$ are equal constant inputs, and the gradient of $\operatorname{LN}$ blows up near that set. 
    So $\operatorname{LN}$ is not globally Lipschitz, and there is no finite $L_{\operatorname{LN}}$ to absorb. 
    The Lipschitz step in part~(j) breaks. 
    A perturbation of size $\mathcal{O}(1/T)$ to the pre normalization activations can be amplified arbitrarily after $\operatorname{LN}$, so we cannot conclude $\Delta_t^{(\ell+1)} = \mathcal{O}(1/T)$. 
    The $\mathcal{O}(1/T)$ guarantee is lost.

    \paragraph{(ii) Implication for fixed length classification}
    The bound $|\hat{p} - \hat{p}'| = \mathcal{O}(1/T)$ is asymptotic. 
    For any fixed finite length $T^{*}$, the gap $\|\mathbf{h}(\mathbf{y}) - \mathbf{h}(\mathbf{y}')\|$ is merely small (at most $C/T^{*}$ for some constant $C$), not zero. 
    Two consequences follow.

    \textbf{For inputs of length $\le T^{*}$:} there are only finitely many strings in $\{0, 1\}^{\le T^{*}}$, and the function $\mathbf{y} \mapsto \mathbb{1}[\mathbf{y} \in \texttt{Odd-1}]$ is well defined on this finite set. 
    By the universal approximation properties of transformers with sufficient width and depth, there exist parameter settings under which the classifier correctly labels every $\mathbf{y}$ with $|\mathbf{y}| \le T^{*}$. 
    Lemma~1 places no obstruction on this.

    \textbf{For unbounded inputs:} Lemma~1 says no single fixed parameter setting recognizes \texttt{Odd-1} uniformly over all lengths. 
    Concretely, fix any parameters $\theta$ and any decision threshold. 
    Since $|\hat{p} - \hat{p}'| \to 0$ as $T \to \infty$, for sufficiently large $T$ the predictions $\hat{p}, \hat{p}'$ on a string and its single bit flip lie within any prescribed margin and in particular both fall on the same side of $0.5$. 
    The transformer then assigns the same label to two strings of opposite parity, contradicting recognition of \texttt{Odd-1}.

    The asymptotic nature of the result identifies where transformers fail (length uniform recognition) versus where they remain capable (any bounded length subset). 
    This contrasts with formal language hierarchies based on finite state machines, which can recognize \texttt{Odd-1} at every length with the same two state automaton.
\end{subquestion}

\end{question}